\chapter{Conclusion}

We solved the two-dimensional system of equations that describe incompressible fluid dynamics.
To do this, we first dealt with the simpler case of linear advection and were able to show that our methods were accurate compared to analytical solutions.
Here, we also implemented Lie-Splitting to reduce computational cost.
We then extended our numerical procedure to respect the incompressible characteristics:
Chorin's projection method allowed us to reduce the dimensionality of the problem by accepting a numerical error.
The pressure was calculated by solving the Poisson equation that could be derived from the incompressibility constraint.
We can show that the error remains bounded. 
To further minimize it, we implemented spectral methods to solve the Poisson equation.
The resulting divergence can be reduced to within machine precision.
Lastly, these methods proved to be suitable for reproducing Rayleigh-B\'{e}nard convection cells.
This was done by certain parametrization and initialization of the model. 
Initial conditions with small perturbations were introduced to trigger the emergence of convective behaviour.
The parameters $\beta= 50.0$ (related to buoyancy) and $\nu = 0.005$ (viscosity) yielded the desired result and are consistent with Rayleigh theory.
Overall, we consider the project to be a successful numerical solution of the incompressible Navier-Stokes problem.


